\subsection{Hahn-Banach 定理（解析形式）}

定理 1（Hahn-Banach）. --- 设 p 是向量空间 E 上的一个次线性函数。设 V 是 E 的一个向量子空间，f 是 V 上的一个线性型，使得对所有 \(y \in V\) ，有 \(f(y) \leqslant p(y)\) 。则存在 E 上的一个线性型 h，它是 f 的扩张，并且使得 \(h(x) \leqslant p(x)\) 对所有 \(x \in E\) 都成立。

数对 \((x, a)\)（满足 \(p(x) \leqslant a\)）所成的集是向量空间 \(E_{1} = E \times R\) 的一个凸子集 P (II, 第 17 页, prop. 19)，而且它显然是一个尖锥。设 \(V_{1}\) 为子空间 \(V \times R\)（\(E_{1}\) 的子空间），并且令 \(g(y, a) = -f(y) + a\)（对每一点 \((y, a) \in V_{1}\) ）。则 g 是 \(V_{1}\) 上由 \(P \cap V_{1}\) 所定义的预序结构下的正线性型；因为若 \((y, a) \in P \cap V_{1}\) ，则 \(a \geqslant p(y) \geqslant f(y)\) ，因此 \(g(y, a) \geqslant 0\) 。其次设 \((x, a) \in E_{1}\) ；我们证明 \((x, a)\) 在由 P 定义的预序下小于 \(V_{1}\) 的某个点。若 \((x', a') \in V_{1}\) ，则 \((x, a) \leqslant (x', a')\) 等价于 \(p(x' - x) \leqslant a' - a\) ；取 \(a' \geqslant p(-x) + a\) ，即见 \((0, a')\) 属于 \(V_{1}\) 且满足要求。于是可以应用 II, 第 21 页, prop. 1；存在 \(E_{1}\) 上的一个线性型 u，它扩张 g 且对由 P 定义的预序是正的。因此 \(u(0, 1) = g(0, 1) = 1\)，且 u 形如 \(u(x, a) = -h(x) + a\) ，其中 h 是 E 上扩张 f 的线性型；此外，对所有 \(x \in E\) 及所有 \(a \geqslant p(x)\) ，有 \(h(x) \leqslant a\) ，因此 \(h(x) \leqslant p(x)\) 。证毕。

推论 1. --- 设 p 是向量空间 E 上的半范数。设 V 是 E 的向量子空间，f 是 V 上的线性型，使得 \(|f(y)| \leqslant p(y)\) 对所有 \(y \in V\) 成立。则存在 E 上定义的线性型 h，它是 f 的扩张，并且使得 \(|h(x)| \leqslant p(x)\) 对所有 \(x \in E\) 成立。

对半范数 \(q\) 与线性型 \(g\) （\(\mathbf{E}\) 上的）而言，关系 \(g \leqslant q\) 与 \(|g| \leqslant q\) 是一回事。推论由 th. 1 得出。

推论 2. --- 设 p 是向量空间 E 上的半范数。给定一点 \(x_{0} \in E\) ，存在 E 上定义的线性型 f，使得 \(f(x_{0}) = p(x_{0})\) 且 \(|f(x)| \leqslant p(x)\) 对所有 \(x \in E\) 成立。

对 \(x_0\) 生成的向量子空间 V 以及 V 上定义的线性型 \(\xi x_0 \mapsto \xi p(x_0)\) 应用 cor. 1。

推论 3. --- 设 V 是赋范空间 E 的向量子空间，f 是 V 上的连续线性型；则存在 E 上定义的连续线性型 h，它扩张 f 且有相同的范数 (GT, X, § 3.2)。

取 \(p(x) = \| f \| \cdot \| x \|\) 应用 cor. 1，即得 \(\| h \| \leqslant \| f \|\) ；但显然 \(\| h \| \geqslant \| f \|\) ，推论得证。

cor. 3 的结论对赋范空间到任一赋范空间中的连续线性映射未必成立 (IV, 第 55 页, exerc. 16, c) 及 V, 第 65 页, exerc. 22)。

